\section{\texorpdfstring{Poziční číselné soustavy}{Pozicni ciselne soustavy}}

% =====================================================================================================================
	\begin{frame}
    \frametitle{Vyjádření polynomem}

		\begin{minipage}[t][0.4\textheight][t]{\textwidth}
      \small
			\boldmath
			\begin{align*}
			H =& A_{Z}(a_{n} a_{n-1} ... a_{1} a_{0}\textcolor{red}{,} a_{-1} a_{-2} ... a_{-m+1} a_{m}) =\\ \\
				=& a_{n} \cdot Z^n + a_{n-1} \cdot Z^{n-1} + ...\\ 
				+& a_{1} \cdot Z^1 + a_{0} \cdot Z^0\textcolor{red}{,} + a_{-1} \cdot Z^{-1} + a_{-2} \cdot Z^{-2} + ...\\ 
				+& a_{-1} \cdot Z^{-m+1} + a_{-m} \cdot Z^{-m}
			\end{align*}
		\end{minipage}
		
		\begin{minipage}[t][0.5\textheight][t]{\textwidth}
			\boldmath
			\begin{itemize}
				\item $H$ ... hodnota čísla, je nezávislá na soustavě, ve které je vyjádřena,		
				
				\item $A_{Z}(a_{n} a_{n-1} ... a_{1} a_{0}, a_{-1} a_{-2} ... a_{-m+1} a_{m})$ ...
					\begin{itemize}
						\item vyjádření čísla v dané soustavě,
						\item $a_{i}$ jsou koeficienty polynomu a vyjadřují číslice,
						\item \textcolor{red}{\textbf{čárka}} reprezentuje řádovou čárku,
						\item \textcolor{red}{\textbf{pozice}} koeficientu určuje násobení příslušné mocniny základu,
						\item koeficienty se záporným indexem náleží neceločíselné části.
					\end{itemize}
			\end{itemize}
		\end{minipage}
		
	\end{frame}

% =====================================================================================================================
	\begin{frame}
    \frametitle{Vyjádření polynomem}
		
		\begin{minipage}[t][0.4\textheight][t]{\textwidth}
      \small
			\boldmath
			\begin{align*}
			H =& A_{Z}(a_{n} a_{n-1} ... a_{1} a_{0}\textcolor{red}{,} a_{-1} a_{-2} ... a_{-m+1} a_{m}) =\\ \\
				=& a_{n} \cdot Z^n + a_{n-1} \cdot Z^{n-1} + ...\\ 
				+& a_{1} \cdot Z^1 + a_{0} \cdot Z^0\textcolor{red}{,} + a_{-1} \cdot Z^{-1} + a_{-2} \cdot Z^{-2} + ...\\ 
				+& a_{-1} \cdot Z^{-m+1} + a_{-m} \cdot Z^{-m}
			\end{align*}
		\end{minipage}
		
		\begin{minipage}[t][0.5\textheight][t]{\textwidth}
			\boldmath
			\begin{itemize}
				\item $Z$ ... základ (báze) soustavy, udává konstantu jejíž násobky mocnin jsou sčítány. Je charakteristikou číselné soustavy:
					\begin{itemize}
						\item $Z=2$ ... binární soustava,
						\item $Z=8$ ... osmičková soustava,
						\item $Z=10$ ... dekadická soustava,
						\item $Z=16$ ... hexadecimální soustava.
					\end{itemize}
				
				\item Koeficienty $a_{i}$ mohou nabývat pouze hodnoty $0$ až $(Z-1)$.
			\end{itemize}
			
		\end{minipage}
	\end{frame}
  
% =====================================================================================================================
	\begin{frame}
    \frametitle{Vyjádření polynomem}
		
		\begin{minipage}[t][0.4\textheight][t]{\textwidth}
    \textbf{Vyjádření pomocí sumy:}
      \small
			\boldmath
			\begin{align*}
      H =& A_{Z}(a_{n} a_{n-1} ... a_{1} a_{0}\textcolor{red}{,} a_{-1} a_{-2} ... a_{-m+1} a_{m}) =\\ \\
				=& \sum^{n}_{-m} a_i\cdot v_i = \sum^{n}_{-m} a_i\cdot Z^i
			\end{align*}
		\end{minipage}
		
		\begin{minipage}[t][0.5\textheight][t]{\textwidth}
			\boldmath
			\begin{itemize}
				\item $i$ ... řád číslice,
				\item $n$ ... nejvyšší řád s nenulovou číslicí,
        \item $-m$ ... nejnižší řád s nenulovou číslicí, pro celá čísla je  $m=0$,
        \item $v_i = Z^i$ ... váha číslice je rovna základu umocněnému příslušným řádem.
			\end{itemize}
			
		\end{minipage}
	\end{frame}
	
% =====================================================================================================================
	\begin{frame}
    \frametitle{Vyjádření čísla v dekadické soustavě}
		
		\begin{minipage}[t][0.6\textheight][t]{\textwidth}
			\boldmath
			\begin{align*}
			(4528)_{10} =& \\ 
				=& 4 \cdot 10^3 + 5 \cdot 10^2 + 2 \cdot 10^1 + 8 \cdot 10^0 =\\ 
				=& 4528\\ 
			(1101001)_{2} =& \\ 
				=& 1 \cdot 2^6 + 1 \cdot 2^5 + 0 \cdot 2^4 + 1 \cdot 2^3 + 0 \cdot 2^2 +\\
				&+ 0 \cdot 2^1 + 1 \cdot 2^0 =\\ 
				=& 105\\ 
			\end{align*}
		\end{minipage}
		
		\begin{minipage}[t][0.4\textheight][t]{\textwidth}
			\boldmath
			\begin{itemize}
				\item Ve všech soustavách jde jen o vyjádření hodnoty - lidé přemýšlí v desítkové soustavě
			\end{itemize}
			
		\end{minipage}
	\end{frame}
	
% =====================================================================================================================
	\begin{frame}
    \frametitle{Vyjádření čísla v dekadické soustavě}
		
		\begin{minipage}[t][0.6\textheight][t]{\textwidth}
			\boldmath
			\begin{align*}
			(527)_{8} =& \\ 
				=& 5 \cdot 8^2 + 2 \cdot 8^1 + 7 \cdot 8^0 =\\ 
				=& 5 \cdot 64 + 2 \cdot 8 + 7 \cdot 1 =\\
				=& 437\\ 
			(1C5)_{16} =& \\ 
				=& 1 \cdot 16^2 + C \cdot 16^1 + 5 \cdot 16^0 =\\
				=& 1 \cdot 256 + 12 \cdot 16 + 5 =\\ 
				=& 453\\ 
			\end{align*}
		\end{minipage}
		
		\begin{minipage}[t][0.4\textheight][t]{\textwidth}
			\boldmath
			\begin{itemize}
				\item V šestnáckové soustavě chybí číslice pro hodnoty 10 až 15. Nahrazujeme je písmeny \textbf{A} až \textbf{F}.
			\end{itemize}
			
		\end{minipage}
	\end{frame}
	
% =====================================================================================================================
	\begin{frame}
    \frametitle{Číslice šestnáctkové soustavy}
		
		\begin{minipage}[c][0.65\textheight][t]{0.3\textwidth}
			\begin{tabular}{| c | c |}
				\hline
				\textbf{Číslice} & \textbf{Hodnota} \\
				\hline
				0 & 0 \\
				\hline
				1 & 1 \\
				\hline
				2 & 2 \\
				\hline
				3 & 3 \\
				\hline
				4 & 4 \\
				\hline
				5 & 5 \\
				\hline
				6 & 6 \\
				\hline
				7 & 7 \\
				\hline
			\end{tabular}
		\end{minipage} % Pozor, bez nové řádky = vše na stejném řádku
		\begin{minipage}[c][0.65\textheight][t]{0.32\textwidth}
			\begin{tabular}{| c | c |}
				\hline
				\textbf{Číslice} & \textbf{Hodnota} \\
				\hline
				8 & 8 \\
				\hline
				9 & 9 \\
				\hline
				A & 10 \\
				\hline
				B & 11 \\
				\hline
				C & 12 \\
				\hline
				D & 13 \\
				\hline
				E & 14 \\
				\hline
				F & 15 \\
				\hline
			\end{tabular}
		\end{minipage} % Pozor, bez nové řádky = vše na stejném řádku
		\begin{minipage}[c][0.65\textheight][t]{0.28\textwidth}
			Výhody:
			\begin{itemize}
				\item zkrácený zápis velkých čísel,
				\item vyjádření stavu mnoha bitů.
			\end{itemize}
		\end{minipage}
		
		\begin{minipage}[t][0.3\textheight][t]{\textwidth}
			Mezi šestnáctkovou a dvojkovou soustavu je snadný přímý převod, který je popsán dále v prezentaci.
		\end{minipage}
	\end{frame}
  
% =====================================================================================================================
	\begin{frame}
    \frametitle{Předpony binárních násobků IEC 60027-2}
		
		\begin{tabular}{ c  c  c  l  l }
				\hline
				\textbf{Násobek} & \textbf{Předpona} & \textbf{Značka} & \textbf{Zdroj názvu} & \textbf{Původ}\\
				\hline
				$2^{10}$ & kibi & Ki & kilobinary: $(2^{10})^{1}$ & kilo: $(10^{3})^{1}$ \\
        $2^{20}$ & mebi & Mi & megabinary: $(2^{10})^{2}$ & mega: $(10^{3})^{2}$ \\
        $2^{30}$ & gibi & Gi & gigabinary: $(2^{10})^{3}$ & giga: $(10^{3})^{3}$ \\
        $2^{40}$ & tebi & Ti & terabinary: $(2^{10})^{4}$ & tera: $(10^{3})^{4}$ \\
        $2^{50}$ & pebi & Pi & petabinary: $(2^{10})^{5}$ & peta: $(10^{3})^{5}$ \\
        $2^{60}$ & exbi & Ei & exabinary: $(2^{10})^{6}$ & exa: $(10^{3})^{6}$ \\
        $2^{70}$ & zebi & Zi & zettabinary: $(2^{10})^{7}$ & zetta: $(10^{3})^{7}$ \\
        $2^{80}$ & yobi & Yi & yottabinary: $(2^{10})^{8}$ & yotta: $(10^{3})^{8}$ \\
				\hline
			\end{tabular}
	\end{frame}